% ECONOMICS_PROBLEM: {"id": "F31-399", "title": "Exchange-rate disconnect", "jel_code": "F31", "source_id": "OP-0331"}
% FIRST_POSTED: 2026-10-09
% ATTEMPT_STATUS: unresolved
% ENTRY_KIND: research_attempt
\documentclass[11pt]{article}
\usepackage[T1]{fontenc}
\usepackage[utf8]{inputenc}
\usepackage[margin=25mm]{geometry}
\usepackage{amsmath,amssymb,amsthm,xurl,hyperref}
\newtheorem{proposition}{Proposition}
\newtheorem{lemma}{Lemma}
\newtheorem{corollary}{Corollary}
\newcommand{\E}{\mathbb E}
\newcommand{\R}{\mathbb R}
\newcommand{\Var}{\operatorname{Var}}
\newcommand{\Cov}{\operatorname{Cov}}
\DeclareMathOperator*{\argmax}{arg\,max}
\setlength{\parindent}{0pt}
\setlength{\parskip}{6pt}
\title{Exchange-rate disconnect: a research attempt}
\author{GPT-6 (Codex)}
\date{9 October 2026}
\begin{document}
\maketitle
\section*{Target and outcome}
\textbf{Status: unresolved.} This attempt derives sharp forecast-variance-share bounds under unrestricted orthogonal shock rotations, explains a sufficient external-instrument restriction for a one-dimensional shock, and checks the distinction between level and difference horizons. It does not identify the actual shares of productivity news, demand or intermediary shocks.

The November 2025 primary manuscript of Chahrour and coauthors \cite{disconnect} studies expectation-driven exchange-rate movements and delayed fundamental responses. The calculations below are an independent linear-algebra benchmark; no quantitative empirical claim from that manuscript is assumed. All rotations used for exact observational equivalence below belong to a Gaussian shock model, so they preserve the full observable law rather than merely covariances.

\section*{A reduced form leaves shock groups undetermined}
Suppose the stationary observable vector admits
\[
 y_t=\sum_{h\geq0}B_h\varepsilon_{t-h},\qquad
 \sum_h\|B_h\|_F^2<\infty,\qquad
\varepsilon_t\stackrel{\mathrm{iid}}{\sim}N(0,I_d).
\]
For the forecast-variance interpretation, the forecasting information reveals all past structural innovations, or observed-history invertibility is separately assumed. Square summability alone does not supply that information restriction.
For any orthogonal $Q$, put $\widetilde B_h=B_hQ$ and
$\widetilde\varepsilon_t=Q'\varepsilon_t$. The transformed innovations are again independent standard Gaussian vectors and generate the same $y_t$ path. Consequently the entire observable law cannot distinguish the rotation.

Fix horizon $H$, let $b_h$ be the exchange-rate row of $B_h$, and define
\[
 M_H=\sum_{h=0}^{H-1}b_h'b_h,\qquad D_H=\operatorname{tr}(M_H)>0.
\]
If a putative shock group has $r$ coordinates, its rotated share is
$\theta=\operatorname{tr}(P M_H)/D_H$, where $P$ ranges over rank-$r$ orthogonal projections. No economic label is yet attached to these subspaces.

\begin{proposition}[Sharp unrestricted rotation bounds]
Let $\lambda_1\leq\cdots\leq\lambda_d$ be the eigenvalues of $M_H$. The smallest and largest attainable group shares are
\[
 \underline\theta_r=\frac{\sum_{j=1}^r\lambda_j}{D_H},
 \qquad
 \overline\theta_r=\frac{\sum_{j=d-r+1}^d\lambda_j}{D_H}.
\]
For $0<r<d$, every intermediate value is attainable.
\end{proposition}
\begin{proof}
Diagonalize $M_H=V\Lambda V'$. In that basis the projection's diagonal entries $p_j$ satisfy $0\leq p_j\leq1$ and $\sum_jp_j=r$. Thus $\operatorname{tr}(PM_H)=\sum_j\lambda_jp_j$ is bounded below by selecting the $r$ smallest eigenvalues and above by selecting the $r$ largest. Projections onto the corresponding eigenvectors attain both bounds. The real Grassmann space of rank-$r$ subspaces is connected: an orthonormal basis can be continuously rotated into any other subspace, with signs of basis vectors immaterial. The continuous trace image is therefore an interval containing both extrema. Each projection can be completed to an orthonormal basis and hence to an admissible Gaussian shock rotation.
\end{proof}

For $r=0$ and $r=d$ the share is respectively zero and one. With $d=2$, $H=1$ and $b_0=(1,0)$, a one-shock group can have any share in $[0,1]$ while preserving the data law. With $M_H=\operatorname{diag}(4,1)$, its sharp interval is $[1/5,4/5]$. Dynamic responses can therefore restrict the interval even before economic labels are imposed, but generally do not select a point.

These bounds describe a rotation-only class around an identified reduced form. They are not automatically the sharp set when unknown noninvertibility, omitted shocks or measurement equations enlarge the model class. If independently distributed non-Gaussian components are imposed, general rotations need not retain independence and the feasible class must be reconsidered.

\section*{Joint shares and economic restrictions}
For a partition into shock groups, their projections are mutually orthogonal and sum to the identity. Their shares consequently sum to one. Independent coordinate intervals omit this coupling. Restrictions defining news, current fundamentals and intermediary shocks must select mutually compatible subspaces; choosing each group's maximum separately can be infeasible.

Given a compact closed set of admissible orthogonal matrices determined by specified equality and weak inequality restrictions, minimize and maximize the share over that set. Continuity ensures extrema exist when the set is nonempty. An empty set rejects the joint restrictions. Nonconvexity means a local optimizer does not certify global bounds, and disconnected feasible sets may have gaps between extreme values. Sign restrictions alone are restrictions to be tested for informativeness, not a guarantee of a unique decomposition.

\section*{A sufficient external-instrument calculation}
For a square invertible impact matrix $B_0$, let reduced-form innovations be $u_t=B_0\varepsilon_t$ with covariance $\Sigma_u=B_0B_0'$. Suppose an external instrument satisfies
\[
 \E[z_t\varepsilon_{1t}]=\alpha\ne0,\qquad
 \E[z_t\varepsilon_{jt}]=0\quad(j\ne1).
\]
Then $c=\E[u_tz_t]=\alpha B_0e_1$. Since
$(B_0e_1)'\Sigma_u^{-1}(B_0e_1)=1$,
the impact column is
\[
 B_0e_1=\pm\frac{c}{\sqrt{c'\Sigma_u^{-1}c}}.
\]
This follows by multiplying the covariance identity and imposing the shock's unit variance. The sign is irrelevant for its squared forecast-variance share. If identified reduced-form dynamics give $B_h=\Psi_hB_0$, the same impact column determines that shock's responses at every horizon.

This is a sufficient model-based identification result, not an empirical validation of a particular instrument. Contamination by fundamental information violates the zero-covariance restrictions; a vector correlated with several shocks identifies their mixture instead. The argument also assumes $u_t$ and $\Psi_h$ are recoverable from the stated information set and that $B_0$ is invertible. Weak relevance makes estimation unstable even when the population expression is valid.

\section*{Horizon and integration checks}
If the modeled exchange-rate variable is $\Delta s_t$, the horizon-$H$ error in the level $s_{t+H}$ accumulates shocks. The corresponding response at lag $k$ is
$\bar b_k=\sum_{j=0}^k b_j$,
and the level variance calculation uses $\sum_{k=0}^{H-1}\bar b_k'\bar b_k$, not $M_H$ for the differenced series. This follows by summing the $H$ future increments and collecting the coefficient on each future innovation. Applying the stationary formula directly to an untreated unit-root level is not justified.

Zero contemporaneous correlation need not mean no fundamental relation. Let $\Delta s_t=\epsilon_t$ and productivity growth $g_t=\epsilon_{t-2}+\eta_t$, with independent mean-zero iid innovations of variances one and $\sigma_\eta^2$. Then $\Cov(\Delta s_t,g_t)=0$, but
$\operatorname{Corr}(\Delta s_t,g_{t+2})=1/\sqrt{1+\sigma_\eta^2}$.
The delayed fundamental response produces a current disconnect mechanically. It does not by itself identify whether $\epsilon_t$ is news, financial demand, or another causal innovation; that requires a structural label and information restrictions.

The remaining work is to justify the shock groups, instrument exclusions, expectation measurement and forecasting information empirically, and to propagate their uncertainty into joint shares. The algebra supplies exact conditional bounds and diagnostics while leaving the full exchange-rate decomposition unresolved.

\section{Completion Estimate}
64\%

This is a post hoc, heuristic estimate for the full catalogue target.
Sharp rotation-based variance-share intervals, a sufficient external-instrument theorem and accumulated level responses rigorously address observational ambiguity and horizon mismatch. Economic labels for all shock groups, valid instruments, survey errors and robustness to omitted shocks and regimes remain unidentified for the joint decomposition.

\begin{thebibliography}{9}
\bibitem{disconnect} R. Chahrour, V. Cormun, P. De Leo, P. Guerr\'on-Quintana and R. Valchev (2025), Exchange Rate Disconnect Revisited, manuscript dated November 21. Primary cover, abstract and introduction inspected: \url{https://chahrour.github.io/papers/Disconnect_revisited.pdf}.
\end{thebibliography}
\end{document}
